% from ega3-1-fr.tex
\oldpage[III]{94}
et comme par définition (0\textsubscript{I}, 4.4.3), $v'^*(\rho)$ est
l'inverse de $\sigma$, on voit que le morphisme composé considéré plus haut
n'est autre finalement que $\theta_q$~; $\psi^\sharp$ est par suite le
$B$-homomorphisme associé
\[
H^q(X,\mathcal F)\otimes_A B\longrightarrow H^q(X',\mathcal F').
\]
Comme $u$ est de type fini, $X$ est réunion finie d'ouverts affines $U_i$
($1\leq i\leq r$)~; soit $\mathfrak U$ le recouvrement $(U_i)$. D'ailleurs
$v$ est un morphisme affine, donc il en est de même de $v'$ (II, 1.6.2,
(iii)), et les $U_i'=v'^{-1}(U_i)$ forment par suite un recouvrement ouvert
affine $\mathfrak U'$ de $X'$. On sait alors (0, 12.1.4.2) que le diagramme
\[
\xymatrix{
H^q(\mathfrak U,\mathcal F)\ar[r]^{\theta_q}\ar[d]
&H^q(\mathfrak U',\mathcal F')\ar[d]\\
H^q(X,\mathcal F)\ar[r]^{\theta_q}&H^q(X',\mathcal F')
}
\]
est commutatif, et les flèches verticales sont des isomorphismes puisque
$X$ et $X'$ sont des schémas (\hyperref[III.1.4.1-fr]{1.4.1}). Il suffit par
suite de prouver que le $\varphi$-morphisme canonique
$\theta_q:H^q(\mathfrak U,\mathcal F)\to H^q(\mathfrak U',\mathcal F')$
est tel que le $B$-homomorphisme associé
\[
H^q(\mathfrak U,\mathcal F)\otimes_A B
\longrightarrow H^q(\mathfrak U',\mathcal F')
\]
soit un isomorphisme. Pour toute suite
$\mathbf s=(i_k)_{0\leq k\leq p}$ de $p+1$ indices de $[1,r]$, posons
\[
U_{\mathbf s}=\bigcap_{k=0}^pU_{i_k},\qquad
U_{\mathbf s}'=\bigcap_{k=0}^pU_{i_k}'=v'^{-1}(U_{\mathbf s}),\qquad
M_{\mathbf s}=\Gamma(U_{\mathbf s},\mathcal F),\qquad
M_{\mathbf s}'=\Gamma(U_{\mathbf s}',\mathcal F').
\]
L'application canonique
$M_{\mathbf s}\otimes_A B\to M_{\mathbf s}'$ est un isomorphisme
(I, 1.6.5), donc l'application canonique
$C^p(\mathfrak U,\mathcal F)\otimes_A B\to
C^p(\mathfrak U',\mathcal F')$ est un isomorphisme, par lequel
$d\otimes1$ s'identifie à l'opérateur cobord
$C^p(\mathfrak U',\mathcal F')\to C^{p+1}(\mathfrak U',\mathcal F')$.
Comme $B$ est un $A$-module \emph{plat}, il résulte aussitôt de la définition
des modules de cohomologie que l'application canonique
$H^q(\mathfrak U,\mathcal F)\otimes_A B\to
H^q(\mathfrak U',\mathcal F')$ est un isomorphisme
(0\textsubscript{I}, 6.1.1). Ce résultat sera généralisé plus tard (\S~6).
\end{proof}

\begin{corollary}[1.4.16]
\label{III.1.4.16-fr}
Soient $A$ un anneau, $X$ un $A$-schéma de type fini, $B$ une $A$-algèbre
fidèlement plate sur $A$. Pour que $X$ soit affine, il faut et il suffit que
$X\otimes_A B$ le soit.
\end{corollary}

\begin{proof}
La condition est évidemment nécessaire (I, 3.2.2)~; montrons qu'elle est
suffisante. Comme $X$ est séparé sur $A$ et que le morphisme
$\operatorname{Spec}(B)\to\operatorname{Spec}(A)$ est plat, il résulte de
(1.4.1) que l'on a
\begin{equation}
\label{III.1.4.16.1-fr}
\tag{1.4.16.1}
H^i(X\otimes_A B,\mathcal F\otimes_A B)
=H^i(X,\mathcal F)\otimes_A B
\end{equation}
pour tout $i\geq0$ et tout $\mathcal O_X$-Module quasi-cohérent $\mathcal F$.
Si $X\otimes_A B$ est affine, le premier membre de
(\hyperref[III.1.4.16.1-fr]{1.4.16.1}) est nul pour $i=1$, donc il en est de
même de $H^1(X,\mathcal F)$ puisque $B$ est un $A$-module fidèlement plat.
Comme $X$ est un schéma quasi-compact, on conclut par le critère de Serre
(II, 5.2.1).
\end{proof}

\begin{proposition}[1.4.17]
\label{III.1.4.17-fr}
Soient $X$ un préschéma,
$0\to\mathcal F\xrightarrow{u}\mathcal G\xrightarrow{v}\mathcal H\to0$
une suite exacte de $\mathcal O_X$-Modules. Si $\mathcal F$ et $\mathcal H$
sont quasi-cohérents, il en est de même de $\mathcal G$.
\end{proposition}

